\chapter{Comunicación RESTful Entre Microservicios}
\label{ch:rest}

\section{Feign Client como mecanismo de comunicación}

La comunicación entre microservicios se implementa de forma declarativa mediante
\textbf{Spring Cloud OpenFeign}, en lugar de construir las llamadas HTTP a mano con
\texttt{RestTemplate} o \texttt{WebClient}. Un \texttt{@FeignClient} se declara como una
interfaz Java; Spring Cloud genera la implementación en tiempo de ejecución y resuelve
la dirección real del servicio destino consultando Eureka -- de modo que ningún
microservicio necesita conocer el host o puerto de otro:

\begin{lstlisting}[language=Java, caption={CatalogoClient.java (inventario) -- Feign client}]
@FeignClient(name = "bookpoint-ms-catalogo")
public interface CatalogoClient {

    @GetMapping("/api/v1/catalogo/{id}")
    EntityModel<LibroRentDTO> obtenerLibroPorId(@PathVariable("id") Long id);
}
\end{lstlisting}

La Tabla~\ref{tab:dependencias-feign} resume las dependencias reales entre
microservicios, todas implementadas mediante este mecanismo.

\begin{table}[h]
\centering
\begin{tabular}{ll}
\toprule
\textbf{Microservicio} & \textbf{Depende vía Feign de} \\
\midrule
Inventario & Catálogo \\
Carrito & Clientes, Catálogo \\
Pedidos & Catálogo, Inventario \\
Pagos & Pedidos \\
Envíos & Pedidos \\
Reseñas & Catálogo, Clientes \\
\bottomrule
\end{tabular}
\caption{Dependencias reales entre microservicios vía Feign}
\label{tab:dependencias-feign}
\end{table}

Los microservicios de sucursales y notificaciones no consumen a ningún otro servicio:
sus dominios son autocontenidos.

\section{Manejo de errores en comunicación remota}

Una llamada Feign puede fallar de dos maneras cualitativamente distintas, y el proyecto
las distingue explícitamente en lugar de tratarlas de forma genérica:

\begin{enumerate}
    \item \textbf{El recurso remoto no existe} (el servicio remoto respondió, pero con
    \texttt{404}): se traduce a una excepción de negocio local
    (\texttt{RecursoNoEncontradoException}), que a su vez el
    \texttt{GlobalExceptionHandler} traduce a un \texttt{404} en la respuesta del
    servicio que hizo la llamada.
    \item \textbf{No se pudo establecer comunicación} (el servicio remoto no responde,
    está caído, o Eureka aún no propagó su registro): se traduce a un error de servidor
    (\texttt{500}), porque no es un problema del dato solicitado sino de la
    infraestructura.
\end{enumerate}

\begin{lstlisting}[language=Java, caption={InventarioServiceImpl.java -- distinción entre 404 remoto y error de red}]
try {
    libroRemoto = catalogoClient.obtenerLibroPorId(inventarioDTO.getLibroId());
} catch (feign.FeignException.NotFound e) {
    throw new RecursoNoEncontradoException(
        "El libro con ID " + inventarioDTO.getLibroId()
        + " no existe en el catalogo.");
} catch (Exception e) {
    throw new RuntimeException(
        "No se pudo conectar con el catalogo. Error de red: "
        + e.getMessage());
}
\end{lstlisting}

Esta distinción no es cosmética: un cliente de la API que recibe un \texttt{404} sabe
que el dato solicitado no existe y puede corregir su petición; uno que recibe un
\texttt{500} sabe que debe reintentar más tarde, porque el problema es de
disponibilidad, no de los datos enviados.

\section{Validación con Postman}

La colección Postman incluye una carpeta dedicada, \textit{"Errores y validaciones
cruzadas (404/400)"}, con casos negativos verificados en vivo contra el stack completo:

\begin{itemize}
    \item Registrar stock de un libro inexistente en inventario $\rightarrow$
    \texttt{404} (validación cruzada contra catálogo).
    \item Crear un pedido con un libro inexistente $\rightarrow$ \texttt{404}
    (validación cruzada contra catálogo).
    \item Crear un pedido con una cantidad mayor al stock disponible
    $\rightarrow$ \texttt{400} ("Stock insuficiente").
    \item Procesar un pago de un pedido inexistente $\rightarrow$ \texttt{404}
    (validación cruzada contra pedidos).
\end{itemize}

Cada uno de estos casos fue ejecutado contra los diez microservicios levantados en
Docker (no simulado) antes de incorporarlo a la colección, incluyendo la verificación
manual de que el mensaje de error y el código HTTP devueltos coincidieran exactamente
con lo documentado en este informe.
